\subsection*{成对数据的相关分析}

\begin{enumerate}
    \item 判断(正确的打 “√”，错误的打“×”).
    \begin{dingautolist}{172}
        \item 变量之间只有函数关系，不存在相关关系. \blkda{×}
        \item 两个变量之间产生相关关系的原因受许多不确定的随机因素的影响. \blkda{√}
        \item 两个变量的相关系数越大，它们的相关程度越强.\blkda{×}
        \item 若相关系数 $r = 0$ ，则两变量 $x$ ， $y$ 之间没有关系. \blkda{×}
    \end{dingautolist}
    \begin{jieda}
        \begin{dingautolist}{172}
            \item 存在相关关系
            \item 相关关系属于非确定性依存关系，其产生受到多种不确定的随机因素影响，不存在严格的对应规则
            \item 相关系数的绝对值越大，相关程度越强
            \item 不存在线性关系，但可以存在其他关系
        \end{dingautolist}
    \end{jieda}
    \item 如果两种证券在一段时间内收益数据的相关系数为正数，则表明\dotfill{(\kh{C})}
    \xx{两种证券之间存在着完全同向的联动关系}{两种证券之间存在着完全反向的联动关系}{两种证券的收益有同向变动的倾向}{两种证券的收益有反向变动的倾向}
    \item 下列各图中所示的两个变量具有相关关系的是\dotfill{(\kh{D})}
    \fourtu{\includegraphics[width=0.4\textwidth]{bodyxb/src/bo_d7m200alb0pc73f2msvg_0_501_998_159_179_0.jpg}}{\includegraphics[width=0.4\textwidth]{bodyxb/src/bo_d7m200alb0pc73f2msvg_0_670_999_152_162_0.jpg}}{\includegraphics[width=0.4\textwidth]{bodyxb/src/bo_d7m200alb0pc73f2msvg_0_838_1002_150_160_0.jpg}}{\includegraphics[width=0.4\textwidth]{bodyxb/src/bo_d7m200alb0pc73f2msvg_0_994_1002_157_176_0.jpg}}
    \xx{(1)(2)}{(1)(3)}{(2)(4)}{(2)(3)}
    \item 从左至右，观察下列三个散点图，变量 $x$ 与 $y$ 的关系依次为\blkda{\ding{172}\ding{174}\ding{173}}(正相关记作 \ding{172} ；负相关记作 \ding{173}；不相关记作 \ding{174}).
    
    \begin{center}
    \includegraphics[width=0.6\textwidth]{bodyxb/src/bo_d7m200alb0pc73f2msvg_0_565_1388_579_169_0.jpg}
    \end{center}
    
    \item 下是收集到的某物品的销售价格 $y$ 和物品的大小 $x$ 的数据:
    
    \begin{center}
    \begin{tabular}{|c|c|c|c|c|c|}
    \hline
    物品大小/$m^2$ & 11.5 & 110 & 80 & 135 & 105 \\
    \cline{1-6}
    销售价格/万元 & 4.8 & 21.6 & 18.4 & 29.2 & 22 \\
    \cline{1-6}
    \hline
    \end{tabular}
    \end{center}

    则根据数据可以判断 $x$ ， $y$ \blkda{有}相关关系. (填“有”或“无”)
    
    \item 对四组数据进行统计, 获得以下散点图, 关于其相关系数比较, 正确的是\dotfill{(\kh{A})}
    
    \begin{center}
    \includegraphics[width=0.6\textwidth]{bodyxb/src/bo_d7m200alb0pc73f2msvg_1_531_270_658_342_0.jpg}
    \end{center}

    \xx{${r}_{2} < {r}_{4} < 0 < {r}_{3} < {r}_{1}$}{${r}_{4} < {r}_{2} < 0 < {r}_{1} < {r}_{3}$}{${r}_{4} < {r}_{2} < 0 < {r}_{3} < {r}_{1}$}{${r}_{2} < {r}_{4} < 0 < {r}_{1} < {r}_{3}$}
    
    \item 已知 5 对数据 $\left( {x,y}\right)$ 的散点图如图,若去掉点 $E\left( {3,{0.5}}\right)$ ,则下列说法正确的是\dotfill{(\kh{B})}
    
    \begin{center}
    \includegraphics[width=0.2\textwidth]{bodyxb/src/bo_d7m200alb0pc73f2msvg_1_234_881_336_292_0.jpg}
    \end{center}
    \xx{变量 $x$ 与变量 $y$ 呈正相关}{变量 $x$ 与变量 $y$ 的相关性变强}{残差平方和变大}{样本相关系数 $r$ 变大}
    \begin{jieda}
        由散点图可知,去掉点 $E\left( {3,{0.5}}\right)$ 后, $y$ 与 $x$ 的线性相关性加强,且为负相关,所以${B}$ 正确, ${A}$ 错误;
        
        由于 $y$ 与 $x$ 的线性相关性加强,所以残差平方和变小,所以 ${C}$ 错误
        
        由于 $y$ 与 $x$ 的线性相关性加强,且为负相关,所以相关系数的绝对值变大，而相关系数为负的,所以样本相关系数 $r$ 变小,所以 ${D}$ 错误. 故选 ${B}$ .
    \end{jieda}
    \item 如图是两个变量的散点图, $y$ 关于 $x$ 的回归方程可能是\dotfill{(\kh{C})}
    \begin{center}
        \includegraphics[width=0.2\textwidth]{bodyxb/src/bo_d7m200alb0pc73f2msvg_1_250_1227_270_343_0.jpg}
    \end{center}

    \xx{$y = 3\ln \left| x\right|  + 2$}{$y = 3{e}^{x} - 1$}{$y =  - 2{x}^{3} + 2$}{$y =  - \dfrac{1}{10}x + 2$}
    \begin{jieda}
        由散点图可知, $y$ 与 $x$ 负相关,故排除 ${A},{B}$ ,对于 ${D} : y =  - \dfrac{1}{10}x + 2$ ,点$(x,y)$偏离$y =  - \dfrac{1}{10}x + 2$ 较大,而点$(x,y)$近似在曲线 $y =  - 2{x}^{3} + 2$ 附近,所以 $y$ 关于 $x$ 的回归方程是 ${C}$的可能性大.
    \end{jieda}
    
    \item 某校地理小组在某座山测得海拔高度、气压和沸点的六组数据绘制成散点图如图，则下列说法不正确的是\dotfill{(\kh{A})}
    
    \twotu{\includegraphics[width=0.55\textwidth]{bodyxb/src/bo_d7m200alb0pc73f2msvg_1_240_1633_357_343_0.jpg}}{\includegraphics[width=0.7\textwidth]{bodyxb/src/bo_d7m200alb0pc73f2msvg_1_596_1653_430_323_0.jpg}}
    \xx{气压与海拔高度呈正相关}{沸点与气压呈正相关}{沸点与海拔高度呈负相关}{沸点与海拔高度、沸点与气压的相关性都很强}
    \begin{jieda}
        由图 1 知气压随海拔高度的增加而减小, 由图 2 知沸点随气压的升高而升高，所以气压与海拔高度呈负相关, 沸点与气压呈正相关, 沸点与海拔高度呈负相关.
        
        由于两个散点图中的点都呈线性分布,所以沸点与海拔高度、沸点与气压的相关性都很强,故 ${B},{C},{D}$ 正确, ${A}$ 错误.故选：A.
    \end{jieda}
    
    \item 已知 ${r}_{1}$ 为随机变量 $X$ 和 $Y$ 的样本相关系数， ${r}_{2}$ 为随机变量 $M$ 和 $N$ 的样本相关系数，则下列说法正确的是\dotfill{(\kh{B})}
    \xx{若 ${r}_{1} > 0$ ，则 $X$ 和 $Y$ 负相关}{若 ${r}_{2} = 0$ ，则 $M$ 和 $N$ 线性不相关}{若 ${r}_{1} = {0.92},{r}_{2} =  - {0.96}$ ，则 $X$ 和 $Y$ 的线性相关程度比 $M$ 和 $N$ 的线性相关程度强}{若 $\left| {r}_{2}\right|$ 越接近 1,则 $M$ 和 $N$ 的线性相关程度越弱}
    \begin{jieda}
        ${A}$ ,若 ${r}_{1} > 0$ ,则 $X$ 和 $Y$ 正相关,故 ${A}$ 错误
                
        B,若 ${r}_{2} = 0$ ,则 $M$ 和 $N$ 线性不相关,故 ${B}$ 正确；
        
        C. 若 ${r}_{1} = {0.92},{r}_{2} =  - {0.96}$ ,则 $\left| {r}_{1}\right|  < \left| {r}_{2}\right|$ ,所以 $X$ 和 $Y$ 的线性相关程度比 $M$ 和 $N$ 的线性相关程度弱,故 ${C}$ 错误;
        
        D. 若 $\left| {r}_{2}\right|$ 越接近 1,则 $M$ 和 $N$ 的线性相关程度越强,故 ${D}$ 错误.
    \end{jieda}
    
    \item 某工厂统计了甲产品在 2024 年 7 月至 12 月的销售量 $y$ (单位: 万件),得到以下数据:
    
    \begin{center}
    \begin{tabular}{|c|c|c|c|c|c|c|}
    \hline
    月份 $x$ & 7 & 8 & 9 & 10 & 11 & 12 \\
    \cline{1-7}
    销售量 $y$ & 11 & 12 & 14 & 15 & 18 & 20 \\
    \cline{1-7}
    \hline
    \end{tabular}
    \end{center}
    
    根据表中所给数据，可得相关系数 $r \approx$ \blkda{$0.99$}. (结果用四舍五入法保留 2 位小数) (参考公式: 相关系数 $r = \dfrac{\mathop{\sum }\limits_{{i = 1}}^{n}\left( {{x}_{i} - \bar{x}}\right) \left( {{y}_{i} - \bar{y}}\right) }{\sqrt{\mathop{\sum }\limits_{{i = 1}}^{n}{\left( {x}_{i} - \bar{x}\right) }^{2}}\sqrt{\mathop{\sum }\limits_{{i = 1}}^{n}{\left( {y}_{i} - \bar{y}\right) }^{2}}}$ ,参考数据: $\mathop{\sum }\limits_{{i = 1}}^{6}\left( {{x}_{i} - \bar{x}}\right) \left( {{y}_{i} - \bar{y}}\right)  = {32}$ , $\sqrt{42} \approx  {6.48})$
    \begin{jieda}
        由已知可得, $\bar{x} = \dfrac{7 + 8 + 9 + {10} + {11} + {12}}{6} = {9.5}$ ,
        $\bar{y} = \dfrac{{11} + {12} + {14} + {15} + {18} + {20}}{6} = {15},$
        则 $\mathop{\sum }\limits_{{i = 1}}^{n}{\left( {x}_{i} - \bar{x}\right) }^{2} = {\left( 7 - {9.5}\right) }^{2} + {\left( 8 - {9.5}\right) }^{2} + {\left( 9 - {9.5}\right) }^{2} + {\left( {10} - {9.5}\right) }^{2} + {\left( {11} - {9.5}\right) }^{2} + {\left( {12} - {9.5}\right) }^{2} = \dfrac{35}{2}$ ,
        $\mathop{\sum }\limits_{{i = 1}}^{n}{\left( {y}_{i} - \bar{y}\right) }^{2} = {\left( {11} - {15}\right) }^{2} + {\left( {12} - {15}\right) }^{2} + {\left( {14} - {15}\right) }^{2} + {\left( {15} - {15}\right) }^{2} + {\left( {18} - {15}\right) }^{2} + {\left( {20} - {15}\right) }^{2} = {60},$
        所以, $r = \dfrac{\mathop{\sum }\limits_{{i = 1}}^{n}\left( {{x}_{i} - \bar{x}}\right) \left( {{y}_{i} - \bar{y}}\right) }{\sqrt{\mathop{\sum }\limits_{{i = 1}}^{n}{\left( {x}_{i} - \bar{x}\right) }^{2}}\sqrt{\mathop{\sum }\limits_{{i = 1}}^{n}{\left( {y}_{i} - \bar{y}\right) }^{2}}} = \dfrac{32}{\sqrt{\dfrac{35}{2}} \times  \sqrt{60}} = \dfrac{32}{5\sqrt{42}} \approx  \dfrac{32}{5 \times  {6.48}} \approx  {0.99}$ .
    \end{jieda}

    \item 据说职工迟到频率与其居住地离上班地点的远近有关. 为验证这个说法, 一位社会学 家作了调查，调查数据见下表.
    
    \begin{center}
    \begin{tabular}{|c|c|c|c|c|c|}
    \hline
    编号 & 年迟到次数 & 住地远近(千米) & 编号 & 年迟到次数 & 住地远近(千米) \\
    \cline{1-6}
    1 & 8 & 1 & 6 & 3 & 10 \\
    \cline{1-6}
    2 & 5 & 3 & 7 & 5 & 12 \\
    \cline{1-6}
    3 & 8 & 4 & 8 & 2 & 14 \\
    \cline{1-6}
    4 & 7 & 6 & 9 & 4 & 14 \\
    \cline{1-6}
    5 & 6 & 8 & 10 & 2 & 18 \\
    \cline{1-6}
    \hline
    \end{tabular}
    \end{center}
    
    则职工年迟到次数与住地远近之间的相关系数为\blkda{$-0.843$}
    \begin{jieda}
        设年迟到次数为 $x$，住地远近（千米）为 $y$，样本量 $n=10$。

数据：
$
(8,1),\ (5,3),\ (8,4),\ (7,6),\ (6,8),\ (3,10),\ (5,12),\ (2,14),\ (4,14),\ (2,18)
$

$
\bar{x} = \dfrac{8+5+8+7+6+3+5+2+4+2}{10} = \dfrac{50}{10} = 5
$

$
\bar{y} = \dfrac{1+3+4+6+8+10+12+14+14+18}{10} = \dfrac{90}{10} = 9
$

$
\begin{aligned}
S_{xy} &= \sum (x_i - \bar{x})(y_i - \bar{y}) \\
&= (8-5)(1-9) + (5-5)(3-9) + (8-5)(4-9) \\
&\quad + (7-5)(6-9) + (6-5)(8-9) + (3-5)(10-9) \\
&\quad + (5-5)(12-9) + (2-5)(14-9) + (4-5)(14-9) \\
&\quad + (2-5)(18-9) \\
&\quad + (-2)\times1 + 0\times3 + (-3)\times5 + (-1)\times5 + (-3)\times9 \\
&= -24 + 0 -15 -6 -1 -2 + 0 -15 -5 -27 \\
&= -95
\end{aligned}
$

$
\begin{aligned}
S_{xx} &= \sum (x_i - \bar{x})^2 \\
&= 3^2 + 0^2 + 3^2 + 2^2 + 1^2 + (-2)^2 + 0^2 + (-3)^2 + (-1)^2 + (-3)^2 \\
&= 9+0+9+4+1+4+0+9+1+9 = 46
\end{aligned}
$

$
\begin{aligned}
S_{yy} &= \sum (y_i - \bar{y})^2 \\
&= (-8)^2 + (-6)^2 + (-5)^2 + (-3)^2 + (-1)^2 + 1^2 + 3^2 + 5^2 + 5^2 + 9^2 \\
&= 64+36+25+9+1+1+9+25+25+81 = 276
\end{aligned}
$

$
r = \dfrac{S_{xy}}{\sqrt{S_{xx}S_{yy}}} = \dfrac{-95}{\sqrt{46 \times 276}}
$

$
46 \times 276 = 12696
$
$
\sqrt{12696} \approx 112.68
$
$
r \approx \dfrac{-95}{112.68} \approx -0.843
$

因此，职工年迟到次数与住地远近之间的相关系数为 $-0.843$.
    \end{jieda}
    
    \item 下表是某国家 18 支足球队职业联赛(比赛采用双循环制，得分计算方法: 每场赛事胜 方得 3 分，负方无分，平局两队各得 1 分)的积分和射门次数，求积分与射门次数的相关系数为\blkda{$0.888$}.
    
    \begin{center}
    \begin{tabular}{|c|c|c|c|c|c|c|c|c|c|}
    \hline
    俱乐部 & A & B & C & D & E & F & G & H & 1 \\
    \cline{1-10}
    积分 & 51 & 64 & 62 & 53 & 47 & 43 & 44 & 42 & 46 \\
    \cline{1-10}
    射门次数 & 418 & 509 & 485 & 425 & 452 & 425 & 393 & 350 & 375 \\
    \cline{1-10}
    \hline
    \end{tabular}
    
    \end{center}
    
    \begin{center}
    \begin{tabular}{|c|c|c|c|c|c|c|c|c|c|}
    \hline
    俱乐部 & J & K & L & M & N & O & P & Q & R \\
    \cline{1-10}
    积分 & 43 & 50 & 35 & 40 & 40 & 32 & 41 & 26 & 32 \\
    \cline{1-10}
    射门次数 & 428 & 415 & 363 & 372 & 377 & 271 & 395 & 306 & 357 \\
    \cline{1-10}
    \hline
    \end{tabular}
    \end{center}
    \begin{jieda}
        $\bar{x} = \dfrac{{51} + {64} + {62} + {53} + {47} + {43} + {44} + {42} + {46} + {43} + {50} + {35} + {40} + {40} + {32} + {41} + {26} + {32}}{18} = \dfrac{791}{18} \approx  {43.94}$
        
$\bar{y} = \dfrac{1}{18}({418} + {509} + {485} + {425} + {452} + {425} + {393} + {350} + {375}$
$+ {428} + {415} + {363} + {372} + {377} + {271} + {395} + {306} + {357})$
$= \dfrac{7116}{18} \approx  {395.33}$

$\mathop{\sum }\limits_{{i = 1}}^{{18}}{x}_{i}{y}_{i} = {51} \times  {418} + {64} \times  {509} + {62} \times  {485} + {53} \times  {425} + {47} \times  {452} + {43} \times  {425} + {44} \times  {393}$
$+ {42} \times  {350} + {46} \times  {375} + {43} \times  {428} + {50} \times  {415} + {35} \times  {363} + {40} \times  {372} + {40} \times  {377} + {32} \times  {271}$
$+ {41} \times  {395} + {26} \times  {306} + {32} \times  {357}$
$= {321316}$

$\mathop{\sum }\limits_{{i = 1}}^{n}{x}_{i}^{2} = {51}^{2} + {64}^{2} + {62}^{2} + {53}^{2} + {47}^{2} + {43}^{2} + {44}^{2} + {42}^{2} + {46}^{2} + {43}^{2} + {50}^{2} + {35}^{2} + {40}^{2} + {40}^{2} + {32}^{2} + {41}^{2} + {26}^{2} + {32}^{2} = {36403}$

$\mathop{\sum }\limits_{{i = 1}}^{{18}}{y}_{i}^{2} = {418}^{2} + {509}^{2} + {485}^{2} + {425}^{2} + {452}^{2} + {425}^{2} + {393}^{2} + {350}^{2} + {375}^{2}$
$+ {428}^{2} + {415}^{2} + {363}^{2} + {372}^{2} + {377}^{2} + {271}^{2} + {395}^{2} + {306}^{2} + {357}^{2} = {2870400}$

$r = \dfrac{\mathop{\sum }\limits_{{i = 1}}^{{18}}{x}_{i}{y}_{i} - {18} \times  \bar{x}\bar{y}}{\sqrt{\mathop{\sum }\limits_{{i = 1}}^{{18}}{x}_{i}^{2} - {18} \times  {\left( \bar{x}\right) }^{2}} \times  \sqrt{\mathop{\sum }\limits_{{i = 1}}^{{18}}{y}_{i}^{2} - {18} \times  {\left( \bar{y}\right) }^{2}}} \\
= \dfrac{{321316} - {18} \times  \dfrac{791}{18} \times  \dfrac{7116}{18}}{\sqrt{{36403} - {18} \times  {\left( \dfrac{791}{18}\right) }^{2}} \times  \sqrt{{2870400} - {18} \times  {\left( \dfrac{7116}{18}\right) }^{2}}} \approx  {0.888}.$
    \end{jieda}

    \item 某基地蔬菜大棚采用无土栽培方式种植各类蔬菜. 过去 50 周的资料显示, 该地周光照量 $X$ (单位:小时)都在 30 小时以上，其中不足 50 小时的有 5 周，不低于 50 小时且不超过 70 小时的有 35 周, 超过 70 小时的有 10 周. 根据统计, 该基地的西红柿增加量 $y$ (千克)与使用某种液体肥料的质量 $x$ (千克)之间的对应数据为如图所示的折线图.
    
    \begin{center}
    \includegraphics[width=0.3\textwidth]{bodyxb/src/bo_d7m200alb0pc73f2msvg_3_249_458_339_250_0.jpg}
    \end{center}
     
    
    依据折线图计算相关系数 $r$ (精确到 0.01 )，并据此判断是否可用线性回归模型拟合 $y$ 与 $x$ 的关系. (若 $\left| x\right|  > {0.75}$ , 对千克则线性相关程度很高, 可用线性回归模型拟合)
    \begin{jieda}
        由已知数据可得 $\bar{x} = \dfrac{2 + 4 + 5 + 6 + 8}{5} = 5, \bar{y} = \dfrac{3 + 4 + 4 + 4 + 5}{5} = 4$
        
        所以相关系数 $r = \dfrac{\mathop{\sum }\limits_{{i = 1}}^{n}\left( {{x}_{i} - \bar{x}}\right) \left( {{y}_{i} - \bar{y}}\right) }{\sqrt{\mathop{\sum }\limits_{{i = 1}}^{n}{\left( {x}_{i} - \bar{x}\right) }^{2}\mathop{\sum }\limits_{{i = 1}}^{n}{\left( {y}_{i} - \bar{y}\right) }^{2}}} = \dfrac{6}{2\sqrt{5} \times  \sqrt{2}} = \sqrt{\dfrac{9}{10}} \approx  {0.95}$
        因为 $r > {0.75}$ ,所以可用线性回归模型拟合 $y$ 与 $x$ 的关系.
    \end{jieda}
    
    \item 国家发改委制定了煤改气、煤改电价格扶持新政策，从而使得煤改气、煤改电用户大幅度增加，下面条形图反映了某省连续 7 个月的煤改气、煤改电的用户数量.
    
    \begin{center}
    \includegraphics[width=0.4\textwidth]{bodyxb/src/bo_d7m200alb0pc73f2msvg_3_251_1326_481_321_0.jpg}
    \end{center}
    
    在给定坐标系中作出煤改气、煤改电用户数量 $y$ 随月份 $t$ 变化的散点图,并用散点图和样本相关系数说明 $y$ 与 $t$ 之间具有线性相关性.
     
    
    \begin{center}
    \includegraphics[width=0.3\textwidth]{bodyxb/src/bo_d7m200alb0pc73f2msvg_3_818_1267_447_408_0.jpg}
    \end{center}
    \begin{jieda}
        图如下：
        \begin{center}
            \includegraphics[width=0.25\linewidth]{bodyxb//src/xianyinghuigui.png}
        \end{center}
        由条形图数据和参考数据得 $\bar{t} = 4,\mathop{\sum }\limits_{{i = 1}}^{7}{\left( {t}_{i} - \bar{t}\right) }^{2} = {28}$ ,又 $\sqrt{\mathop{\sum }\limits_{{i = 1}}^{7}{\left( {y}_{i} - \bar{y}\right) }^{2}} \approx  {0.53},\mathop{\sum }\limits_{{i = 1}}^{7}\left( {{t}_{i} - \bar{t}}\right) \left( {{y}_{i} - \bar{y}}\right)  = {2.79}$ ,所以
        $r = \dfrac{\mathop{\sum }\limits_{{i = 1}}^{7}\left( {{t}_{i} - \bar{t}}\right) \left( {{y}_{i} - \bar{y}}\right) }{\sqrt{\mathop{\sum }\limits_{{i = 1}}^{7}{\left( {t}_{i} - \bar{t}\right) }^{2}\mathop{\sum }\limits_{{i = 1}}^{7}{\left( {y}_{i} - \bar{y}\right) }^{2}}} \approx  \dfrac{2.79}{{0.53} \times  2 \times  {2.646}} \approx  {0.995},$
        因为 $y$ 与 $t$ 的相关系数接近1,说明 $y$ 与 $t$ 的线性相关性很强,从而可以用线性回归模型拟合 $y$ 与 $t$ 的关系.
    \end{jieda}
\end{enumerate}






\clearpage
\subsection*{一元线性回归分析}

【知识点】回归系数 $\hat{a}$ 与 $\hat{b}$ 的计算方法如下：

$$
\begin{cases}
\hat{a} = \frac{\sum\limits_{i=1}^{n}(x_i - \bar{x})(y_i - \bar{y})}{\sum\limits_{i=1}^{n}(x_i - \bar{x})^2} = \frac{\sum\limits_{i=1}^{n}x_i y_i - n\bar{x}\bar{y}}{\sum\limits_{i=1}^{n}x_i^2 - n\bar{x}^2} \\
\\
\hat{b} = \bar{y} - \hat{a}\bar{x} = \frac{\sum\limits_{i=1}^{n}y_i - \hat{a}\sum\limits_{i=1}^{n}x_i}{n}
\end{cases}
$$

其中，$\bar{x}$ 与 $\bar{y}$ 分别是数据 $x_i$ 与 $y_i$ （$i = 1, 2, \cdots, n$）的算术平均数；数对 $(\bar{x}, \bar{y})$
称为样本点的中心。

\begin{enumerate}
\item 有关线性回归的说法，不正确的是\dotfill{(\kh{D})}
\xx{相关关系的两个变量不一定是因果关系}{散点图能直观反映数据的相关程度}{回归直线最能代表线性相关的两个变量之间的关系}{任意一组数据都有回归方程}

\item 下列有关回归分析的说法正确的是\dotfill{(\kh{D})}
\xx{样本相关系数越大，则两变量的相关性就越强}{回归直线就是散点图中经过样本数据点最多的那条直线}{回归直线方程不一定过样本中心点 $(\overline{x},\overline{y})$}{回归分析中，样本相关系数 $r<0$，则两变量是负相关关系}

\item 设一条经验回归直线的方程为 $\hat{y}=3-1.2x$，则当变量x增加一个单位时\dotfill{(\kh{C})}
\xx{y平均增加1.2个单位}{y平均增加3个单位}{y平均减少1.2个单位}{y平均减少3个单位}

\item ①事件A和事件B互斥，则 $P(A\cup B)=P(A)+P(B)$；\\
②数据2, 3, 6, 7, 8, 10, 13, 15的第50百分位数为7；\\
③在线性回归模型中，拟合误差越小，表示回归的效果越好；\\
④随机变量X的方差 $D(X)=1$，则 $D(2X+1)=3$。\\
其中正确命题的序号为\blkda{（1）（3）}

\item 由数据 $(x_{1},y_{1}), (x_{2},y_{2}), \cdots, (x_{10},y_{10})$ 可得y关于x的回归直线方程为 $\hat{y}=-2x+7.2$，若 $\sum_{i=1}^{10}x_{i}=6$，则 $\sum_{i=1}^{10}y_{i}=$\blkda{60}

\item 已知y与x之间的线性回归方程为 $\hat{y}=1.6x+21$，其样本点的中心为 $(\overline{x},37)$，样本数据中的x取值依次为2, 8, 6, 14, m，则$m=$\blkda{20}

\item 根据样本数据，得到回归直线方程 $\hat{y}=\hat{b}x+\hat{a}$，则\dotfill{(\kh{C})}
\begin{center}
\begin{tabular}{|c|c|c|c|c|c|c|}
\hline
$x$ & 3 & 4 & 5 & 6 & 7 & 8 \\
\hline
$y$ & -3.0 & -2.0 & 0.5 & -0.5 & 2.5 & 4.0 \\
\hline
\end{tabular}
\end{center}
\xx{$\hat{a}>0, \hat{b}>0$}{$\hat{a}>0, \hat{b}<0$}{$\hat{a}<0, \hat{b}>0$}{$\hat{a}<0, \hat{b}<0$}

\item 蟋蟀鸣叫可以说是大自然优美、和谐的音乐，殊不知蟋蟀鸣叫的频率x(每分钟鸣叫的次数)与气温y(单位：$^\circ\text{C}$)存在着较强的线性相关关系，某地观测人员根据如表的观测数据，建立了y关于x的线性回归方程 $\hat{y}=0.25x+k$，则下列说法不正确的是\dotfill{(\kh{D})}
\begin{center}
\begin{tabular}{|c|c|c|c|c|c|}
\hline
$x$(次数/分钟) & 20 & 30 & 40 & 50 & 60 \\
\hline
$y(^\circ\text{C})$ & 25 & 27.5 & 29 & 32.5 & 36 \\
\hline
\end{tabular}
\end{center}
\xx{k的值是20}{变量x,y呈正相关关系}{若x的值增加1，则y的值约增加0.25}{当蟋蟀52次/分鸣叫时，该地当时的气温预报值为 $33.5^\circ\text{C}$}

\item 为了更好地适应市场需求，某企业根据市场调研得到研发投入x(亿元)与产品收益y(亿元)的数据统计如下：
\begin{center}
\begin{tabular}{|c|c|c|c|c|c|c|c|}
\hline
$x$ & 1 & 2 & 3 & 4 & 5 & 6 & 7 \\
\hline
$y$ & 2 & 3 & 5 & 7 & 8 & 8 & 9 \\
\hline
\end{tabular}
\end{center}
参考公式：$\hat{b} = \dfrac{\sum_{i=1}^{n}x_{i}y_{i}-n\overline{x}\cdot\overline{y}}{\sum_{i=1}^{n}x_{i}^{2}-n\overline{x}^{2}} = \dfrac{\sum_{i=1}^{n}(x_{i}-\overline{x})(y_{i}-\overline{y})}{\sum_{i=1}^{n}(x_{i}-\overline{x})^{2}}, \hat{a}=\overline{y}-\hat{b}\overline{x}$。已知 $\sum_{i=1}^{n}(x_{i}-\overline{x})(y_{i}-\overline{y})=34, \sum_{i=1}^{n}(x_{i}-\overline{x})^{2}=28$，则下列选项不正确的是\dotfill{(\kh{D})}
\xx{$\overline{x}=4, \overline{y}=6$}{由散点图知变量x和y正相关}{相关系数的绝对值越接近0，表示x,y的线性相关程度越弱}{用最小二乘法求得y关于x的线性回归直线方程为 $\hat{y}=1.5x+0.5$}

\item 给出下列说法：\\
①回归直线 $\hat{y}=\hat{b}x+\hat{a}$ 恒过样本点的中心 $(\overline{x},\overline{y})$，且至少过一个样本点；\\
②两个变量相关性越强，则相关系数就越接近1；\\
③将一组数据的每个数据都加一个相同的常数后，方差不变；\\
④在回归直线方程 $\hat{y}=2-0.5x$ 中，当解释变量x增加一个单位时，预报变量平均减少0.5个单位。\\
其中说法正确的是\blkda{(3)(4)}

\item 若变量x,y线性相关，由数据 $(x_{i},y_{i})(i=1,2,3,4,5)$ 求得回归方程为 $\hat{y}=-2x+5$，则下列结论一定成立的是\dotfill{(\kh{D})}
\xx{$\sum_{i=1}^{5}y_{i}=-2\sum_{i=1}^{5}x_{i}$}{$\sum_{i=1}^{5}y_{i}=-2\sum_{i=1}^{5}x_{i}+5$}{$\sum_{i=1}^{5}y_{i}=-10\sum_{i=1}^{5}x_{i}+25$}{$\sum_{i=1}^{5}y_{i}=-2\sum_{i=1}^{5}x_{i}+25$}

\item 用模型 $y=a\cdot e^{kx}$ 拟合一组数据，令 $z=\ln y$，将模型转化为经验回归方程 $z=0.1x+3$，则 $a\cdot k=$\blkda{$\dfrac{e^3}{10}$}

\item 某高中学校高二五班5名学生的数学和物理成绩如表：
\begin{center}
\begin{tabular}{|c|c|c|c|c|c|}
\hline
学生 & A & B & C & D & E \\
\hline
数学成绩(x) & 88 & 76 & 73 & 66 & 63 \\
\hline
物理成绩(y) & 78 & 65 & 71 & 64 & 61 \\
\hline
\end{tabular}
\end{center}
则物理成绩y对数学成绩x的经验回归方程为\blkda[4]{$y=0.625x+22.06$}

\item 某工厂为研究某种产品的产量(吨)与所需某种原材料的质量(吨)的相关性，在生产过程中收集4组对应数据(x,y)，如下表所示。(离差=观测值-预测值)
\begin{center}
\begin{tabular}{|c|c|c|c|c|}
\hline
$x$ & 3 & 4 & 5 & 6 \\
\hline
$y$ & 2.5 & $m$ & 4 & 4.5 \\
\hline
\end{tabular}
\end{center}
根据表中数据，得出y关于x的经验回归方程为 $\hat{y}=0.7x+\hat{a}$，据此计算出在样本 $(6,4.5)$ 处的离差为-0.05，则表中$m$的值为\blkda{3}

\item 为了解大学校园附近餐馆的月营业收入(单位：千元)和该店周围学校的学生人数(单位：千人)的关系，现抽取了10所大学的记录数据，见下表。
\begin{center}
\begin{tabular}{|c|c|c|c|c|c|c|c|c|c|c|}
\hline
学生人数(千人) & 2 & 6 & 8 & 8 & 12 & 16 & 20 & 20 & 22 & 26 \\
\hline
月营业收入y(千元) & 58 & 105 & 88 & 118 & 117 & 137 & 157 & 169 & 149 & 202 \\
\hline
\end{tabular}
\end{center}
\begin{enumerate}
\item 根据以上资料，建立月营业收入y对该店周围学校的学生人数x的回归方程；
\item 已知某餐馆周围学校的学生人数为10000人，试对该店月营业收入作出预测。
\end{enumerate}
\begin{jieda}
    \begin{enumerate}
        \item $y=5x+60$
        \item 110000元
    \end{enumerate}
\end{jieda}

\item 某种机械设备随着使用年限的增加，它的使用功能逐渐减退，使用价值逐年减少，通常把它使用价值逐年减少的“量”换算成费用，称之为“失效费”。某种机械设备的使用年限x(单位：年)与失效费y(单位：万元)的统计数据如下表所示：
\begin{center}
\begin{tabular}{|c|c|c|c|c|c|c|c|}
\hline
使用年限x(单位：年) & 1 & 2 & 3 & 4 & 5 & 6 & 7 \\
\hline
失效费y(单位：万元) & 2.90 & 3.30 & 3.60 & 4.40 & 4.80 & 5.20 & 5.90 \\
\hline
\end{tabular}
\end{center}
\begin{enumerate}
\item 由上表数据可知，可用线性回归模型拟合y与x的关系，请用相关系数加以说明；(精确到0.01)
\item 求出y关于x的线性回归方程，并估算该种机械设备使用10年的失效费。
\end{enumerate}
\begin{jieda}
    \begin{enumerate}
        \item r=0.99
        \item $y=0.5x+2.3$ ; 7.3万元
    \end{enumerate}
\end{jieda}


    
\end{enumerate}


\clearpage

\subsection*{2×2 列联表}
\begin{enumerate}



\item 下面是 $2\times 2$ 列联表：
\begin{center}
\begin{tabular}{|c|c|c|c|}
\hline
 & 1 & 2 & 合计 \\
\hline
$X_{1}$ & $a$ & 21 & 73 \\
\hline
$X_{2}$ & 22 & 25 & 47 \\
\hline
合计 & $b$ & 46 & 120 \\
\hline
\end{tabular}
\end{center}
则表中$a, b$的值分别为\dotfill{(\kh{C})}
\xx{94, 72}{52, 50}{52, 74}{74, 52}

\item 第31届世界大学生运动会将于2023年7月28日至8月8日在成都举行，组委会安排100名志愿者担任对外翻译工作，在下面“性别与会法语”的 $2\times 2$ 列联表中，$a+b+d=$\blkda{88}
\begin{center}
\begin{tabular}{|c|c|c|c|}
\hline
 & 会法语 & 不会法语 & 总计 \\
\hline
男 & $a$ & $b$ & 40 \\
\hline
女 & 12 & $d$ & \\
\hline
总计 & 36 & & 100 \\
\hline
\end{tabular}
\end{center}

\item 根据分类变量与的观测数据，计算得到 $\chi^{2}\approx 0.837$，依据小概率值 $\alpha=0.1$ $(\chi^{2}_{0.1}=2.706)$ 的独立性检验，则\dotfill{(\kh{B})}
\xx{变量与不独立}{变量与独立}{变量与不独立，这个结论犯错误的概率不超过0.1}{变量与独立，这个结论犯错误的概率不超过0.1}

\item 我校为了研究学生的作业量与玩手机(有关和无关)的关系，运用 $2\times 2$ 列联表进行独立性研究，经计算 $K^{2}=5.059$，则至少有\blkda{$97.5 \%$}的把握认为“作业量与玩手机有关”。
\begin{center}
\begin{tabular}{|c|c|c|c|}
\hline
 & 认为作业多 & 认为作业不多 & 总数 \\
\hline
喜欢玩手机 & 18 & 9 & 27 \\
\hline
不喜欢玩手机 & 8 & 15 & 23 \\
\hline
总数 & 26 & 24 & 50 \\
\hline
\end{tabular}
\end{center}

\begin{center}
\begin{tabular}{|c|c|c|c|c|c|c|c|c|c|c|}
\hline
$P(K^{2}\ge k)$ & 0.50 & 0.40 & 0.25 & 0.15 & 0.10 & 0.05 & 0.025 & 0.010 & 0.005 & 0.001 \\
\hline
$k$ & 0.455 & 0.708 & 1.323 & 2.072 & 2.706 & 3.841 & 5.024 & 6.635 & 7.879 & 10.828 \\
\hline
\end{tabular}
\end{center}

\item 在性别与吃零食这两个分类变量的计算中，下列说法正确的是\dotfill{(\kh{D})}\\
①若的观测值为$k=6.635$，我们有99\%的把握认为吃零食与性别有关系，那么在100个吃零食的人中必有99人是女性；\\
②从独立性检验可知有99\%的把握认为吃零食与性别有关系时，我们说某人吃零食，那么此人是女性的可能性为99\%；\\
③若从统计量中求出有99\%的把握认为吃零食与性别有关系，是指有1\%的可能性使得出的判断出现错误。
\xx{①②}{①③}{②③}{③}

\item 为了验证牛的毛色(黑色、红色)和角(有角、无角)这两对相对性状是否相关，某学院进行了一次数据统计，根据形成的 $2\times 2$ 列联表，计算得到 $\chi^{2}\approx 2.786$，根据小概率值 $\alpha=0.05$ 的独立性检验(已知独立性检验中 $\chi^{2}_{0.05}=3.841$)，下列结论正确的是\dotfill{(\kh{A})}
\xx{牛的毛色与角无关}{牛的毛色与角无关，此推断犯错误的概率不超过0.05}{牛的毛色与角有关}{牛的毛色与角有关，此推断犯错误的概率不超过0.05}

\item 某单位对员工是否愿意被外派与年龄的关系进行了一次调查，根据独立性检验原理，处理所得数据之后发现，得到“是否愿意被外派与年龄有关”这个结论犯错误的概率大于0.001，而不大于0.01，则 $\chi^{2}$ 的值可能为\dotfill{(\kh{C})}\\
附表：
\begin{center}
\begin{tabular}{|c|c|c|c|}
\hline
$P(\chi^{2}\ge\chi_{\alpha})$ & 0.05 & 0.01 & 0.001 \\
\hline
$\chi_{\alpha}$ & 3.841 & 6.635 & 10.828 \\
\hline
\end{tabular}
\end{center}
\xx{3.206}{6.561}{7.879}{11.028}


\item 2025年春节期间，电影《哪吒之魔童闹海》掀起全民观影热潮，连续7天票房逆势攀升，单日最高突破8.6亿元，吸引部分家庭扶老携幼共赴影院，缔造中国影史春节档票房与观影人次双冠王的奇迹。某电影院为了解民众观影的喜欢程度，随机采访了180名观影人员，得到下表：
\begin{center}
\begin{tabular}{|c|c|c|c|}
\hline
 & 喜欢 & 不喜欢 & 合计 \\
\hline
未成年人 & $m$ & 80 & 100 \\
\hline
成年人 & 20 & 80 & $s$ \\
\hline
合计 & $n$ & 180 & $t$ \\
\hline
\end{tabular}
\end{center}
依据小概率值 $\alpha=0.1$ 的独立性检验，能否认为喜欢电影《哪吒之魔童闹海》与是否成年有关？\\
参考公式：$\chi^{2}=\dfrac{n(ad-bc)^{2}}{(a+b)(c+d)(a+c)(b+d)}$
\begin{center}
\begin{tabular}{|c|c|c|c|}
\hline
$\alpha$ & 0.1 & 0.05 & 0.01 \\
\hline
$\chi_{\alpha}$ & 2.706 & 3.841 & 6.635 \\
\hline
\end{tabular}
\end{center}
\begin{jieda}
    $\chi^2\approx 0.643< 3.841$,接受原假设，即认为无关.
\end{jieda}

\item 气管炎是一种常见的呼吸道疾病，医药研究人员对甲、乙两种中草药治疗慢性气管炎的疗效进行了对比，所得数据如下表所示，问：它们的疗效有无显著差异？
\begin{center}
\begin{tabular}{|c|c|c|c|}
\hline
 & 有效 & 无效 & 总计 \\
\hline
甲药 & 184 & 61 & 245 \\
\hline
乙药 & 91 & 9 & 100 \\
\hline
总计 & 275 & 70 & 345 \\
\hline
\end{tabular}
\end{center}
\begin{jieda}
    $\chi^2\approx 11.098> 3.841$,拒绝原假设，即认为疗效有显著差异.
\end{jieda}

\item 某工人在操作方法改进前后生产的零件情况如下：
\begin{center}
\begin{tabular}{|c|c|c|c|}
\hline
 & 合格数 & 不合格数 & 总计 \\
\hline
改进前 & 2422 & 439 & 2861 \\
\hline
改进后 & 2892 & 447 & 3339 \\
\hline
总计 & 5314 & 886 & 6200 \\
\hline
\end{tabular}
\end{center}
问：改进操作方法能否显著降低不合格率？
\begin{jieda}
    $\chi^2\approx 40.26< 3.841$,拒绝原假设，即认为改进操作方法能显著降低不合格率.
\end{jieda}


\item 微信是现代生活进行信息交流的重要工具，据统计，某公司200名员工中90\%的人使用微信，其中每天使用微信时间在一小时以内的有60人，其余每天使用微信在一小时以上，若将员工年龄分成青年(年龄小于40岁)和中年(年龄不小于40岁)两个阶段，使用微信的人中75\%是青年人，若规定：每天使用微信时间在一小时以上为经常使用微信，经常使用微信的员工中 $\dfrac{2}{3}$ 是青年人。
\begin{enumerate}
\item 若要调查该公司使用微信的员工经常使用微信与年龄的关系，列出 $2\times 2$ 列联表；
\begin{center}
\begin{tabular}{|c|c|c|c|}
\hline
 & 青年人 & 中年人 & 合计 \\
\hline
经常使用微信 & & & \\
\hline
不经常使用微信 & & & \\
\hline
合计 & & & \\
\hline
\end{tabular}
\end{center}
\item 由列联表中所得数据，是否有99.9\%的把握认为“经常使用微信与年龄有关”？
\begin{center}
\begin{tabular}{|c|c|c|}
\hline
$P(K^{2}\ge k)$ & 0.010 & 0.001 \\
\hline
$k$ & 6.635 & 10.828 \\
\hline
\end{tabular}
\end{center}
附：$K^{2}=\dfrac{n(ad-bc)^{2}}{(a+b)(c+d)(a+c)(b+d)}$
\end{enumerate}
\begin{jieda}
    \begin{enumerate}
        \item \begin{center}
\begin{tabular}{|c|c|c|c|}
\hline
 & 青年人 & 中年人 & 合计 \\
\hline
经常使用微信 & 80&40 &120 \\
\hline
不经常使用微信 &55 & 5&60 \\
\hline
合计 & 135 &45 &180 \\
\hline
\end{tabular}
\end{center}

\item  $\chi^2\approx 13.333> 10.83$,拒绝原假设,有$99.9\%$的把握认为有关.
    \end{enumerate}
\end{jieda}

\end{enumerate}